\section{Ablation studies}
\label{sec:ablations}

\cond{EBG} combines four things: evidence-basis labels shown to the agent, a reversible notice at the terminal
action, \emph{gating} that enforces only strong-strength resolutions, \emph{safe submission}, and
search-before-submit. The pre-registered exploratory contrasts separate them (\Cref{tab:ablations}). Two
components carry essentially the whole effect; one is inert on success but changes the failure profile; and the
gate --- the component the paper is named for --- turns out to be a deliberate trade rather than an improvement.

\begin{table}[t]
\centering
\caption{Component ablations, all pre-registered as \emph{exploratory} and reported with Benjamini--Hochberg
adjustment. A positive success difference favours the first-named condition; a negative harm difference favours
it. The two bold rows are the results that cut against the proposed design.}
\label{tab:ablations}
\small
\setlength{\tabcolsep}{5pt}
\begin{tabular}{@{}lp{0.29\linewidth}rcrr@{}}
\toprule
\textbf{Contrast} & \textbf{What it isolates} & \textbf{Success} & \textbf{95\% CI} & $\boldsymbol{p_{\text{BH}}}$ & \textbf{Harm} \\
 & & \textbf{(pp)} & & & \textbf{(pp)} \\
\midrule
\cond{EB} $-$ \cond{G}   & evidence-basis vs metadata labels, no enforcement & $+11.8$ & \ci{8.3}{15.2} & $2.2\!\times\!10^{-5}$ & $-5.2$ \\
\cond{EB} $-$ \cond{T}   & evidence-basis vs temporal labels, no enforcement & $+8.5$ & \ci{5.4}{11.6} & $2.2\!\times\!10^{-5}$ & $-4.7$ \\
\cond{EBG} $-$ \cond{EB} & adding gated enforcement to the labels & $+6.0$ & \ci{3.9}{8.2} & $2.2\!\times\!10^{-5}$ & $+0.5$ \\
\cond{EBG} $-$ \cond{R}  & whole mechanism vs prior record-based enforcement & $+16.3$ & \ci{11.7}{21.0} & $2.2\!\times\!10^{-5}$ & $-12.9$ \\
\cond{EBG} $-$ \cond{EBG-noR3} & search-before-submit\footnotemark[1] & $+1.3$ & --- & --- & $-0.3$ \\
\cond{EBG} $-$ \cond{EBG-unsafe} & safe submission & $-0.2$ & \ci{-1.0}{0.7} & $0.776$ & $-0.2$ \\
\midrule
\textbf{\cond{EBG} $-$ \cond{EBG-all}} & \textbf{the gate itself} & $\boldsymbol{-9.7}$ & \textbf{\ci{-15.2}{-4.3}} & $\boldsymbol{0.0013}$ & $\boldsymbol{-3.7}$ \\
\textbf{\cond{LG} $-$ \cond{LA}} & \textbf{gating an LLM arbiter} & $\boldsymbol{-2.4}$ & \textbf{\ci{-5.6}{0.6}} & $\boldsymbol{0.167}$ & $\boldsymbol{+8.1}$ \\
\bottomrule
\end{tabular}
\footnotetext[1]{A descriptive difference of pooled rates (61.3\% against 60.0\%) rather than one of the
pre-registered contrast set, so no interval or adjusted $p$ is computed for it.}

\vspace{2pt}
{\footnotesize Source: \srcfile{PHASE3/ANALYSIS/CONFIRMATORY_primary_cells.json}; uncorrected $p$ for
\cond{EBG} $-$ \cond{EBG-all} is $6.9\times10^{-4}$ and for \cond{LG} $-$ \cond{LA} is $0.124$.}
\end{table}


\subsection{The labels do most of the work}

Replacing metadata labels with evidence-basis labels, with \emph{no enforcement at all}, raises success by
$11.8$\pp{} over \cond{G} and $8.5$\pp{} over the temporal arbiter \cond{T}, and cuts the harmful rate from 7.9\%
and 7.4\% to 2.7\%. False resolution falls from 9.2\% (\cond{G}) and 3.6\% (\cond{T}) to 2.0\%. Most of the
mechanism's benefit is in \emph{what the agent is told}, before anything is enforced. This is the part of the
result we consider best supported, and it is the part with the least novel machinery behind it.

\subsection{Gating adds a real but smaller increment}

Adding the gated notice to the labels adds $6.0$\pp{} \ci{3.9}{8.2} of success with no detectable change in harm
($+0.5$\pp{} \ci{-0.1}{1.2}), lifts conflict preservation from 81.6\% to 88.5\%, and halves step-limit failures.
The notice mostly recovers runs in which the agent had the right labels and still submitted the superseded value.

\subsection{Safe submission buys format robustness, not success}

Success is statistically indistinguishable with and without safe submission ($-0.2$\pp{} \ci{-1.0}{0.7},
$p_{\text{BH}}=0.78$), and so is harm. What changes is the failure profile: format errors nearly double without it,
from 3.3\% to 6.4\% pooled, while unnecessary abstention falls from 12.8\% to 9.9\%. Without the fallback a
mangled reply to a notice becomes a format failure; with it, the pre-notice answer stands, which sometimes means
an abstention. We would keep the component for robustness, but no success claim rests on it.
Search-before-submit is worth about $1.3$\pp{} and reduces unnecessary abstention from 17.6\% to 12.8\%.

\subsection{The gate: the central trade}

Enforcing \emph{every} arbitration rather than only strong-strength ones scores \textbf{$9.7$\pp{} higher} than
\cond{EBG} (71.1\% against 61.3\%; CI \ci{-15.2}{-4.3} in \cond{EBG}'s disfavour, $p_{\text{BH}}=0.0013$). It also

\begin{itemize}[leftmargin=1.4em,itemsep=1pt,topsep=2pt]
\item more than doubles the harmful rate, 3.1\% $\rightarrow$ 6.9\% ($-3.7$\pp{} \ci{-6.6}{-1.3} in
\cond{EBG}'s favour);
\item restores arbitrary picks from 13 runs to 166;
\item raises false resolution from 1.9\% to 12.9\% and drops conflict preservation from 88.5\% to 77.3\%.
\end{itemize}

This is the clearest answer in the study to the second half of our research question, and it goes against the
mechanism we proposed. \textbf{Indiscriminate enforcement does increase harmful answers} --- and on this benchmark
it also produces more correct answers than gated enforcement does. Gating is a trade, not a dominance.

Which side of that trade one should take depends on the relative cost of a wrong confident answer and a missed
correct one, which this benchmark does not price. We state the trade rather than resolve it, with two
qualifications. It is sensitive to the item set: $-4.8$\pp{} \ci{-9.7}{0.3} and not significant on the
double-agreement subset, $-17.1$\pp{} on the flagged complement, which suggests part of it lives in items where
the annotators themselves disagreed. And success and harm are not commensurable: $9.7$\pp{} of success against
$3.8$\pp{} of additional harm and a thirteen-fold increase in arbitrary picks is a judgement, not an arithmetic
comparison.

\subsection{Gating helps only when the arbiter is reliable}

The same gating machinery driven by an LLM arbiter instead of the deterministic one does not help: success is
statistically unchanged ($-2.4$\pp{} \ci{-5.6}{0.6}, $p_{\text{BH}}=0.17$) and harm rises by $8.1$\pp{}
\ci{6.1}{10.2}. Under \cond{LG}, false resolution reaches 37.0\% and conflict preservation falls to 42.4\% --- the
worst of any condition, worse than no mechanism at all. For Qwen2.5-7B specifically, \cond{LG} falsely resolves
75.9\% of genuine conflicts.

The useful generalisation is therefore neither ``enforcement helps'' nor ``gating helps'', but that
\textbf{enforcement inherits the reliability of whatever supplies its labels}. Enforcing a deterministic,
evidence-grounded arbiter's strong resolutions helps; enforcing a record-metadata scanner's (\cond{RT}) or a
language model's (\cond{LG}) is actively harmful. The value of the evidence-basis idea is that it gives
enforcement something defensible to enforce, and a principled place to stop.
